
\chapter{Сверка: где теряются деньги}\label{ch:reconciliation}

Деньги почти никогда не \enquote{теряются} в~одной точке.
Они растворяются в~переходах между системами, каждая из которых
честно хранит свою версию события. В~журнале авторизаций покупка уже
разрешена, в~клиринговом файле она ещё не появилась, а~в~банковской
выписке от неё осталась только нетто-сумма после комиссий и~взаимных
вычетов. Сверка собирает из этих разрозненных следов одну операцию.

Внутренний реестр обязательств из главы~\ref{ch:internal-ledger}
отвечает на вопрос, какое денежное состояние видит сам продукт.
Сверка задаёт следующий вопрос: совпадает ли эта внутренняя картина
с~внешними файлами схемы, эквайера и~банка. Перепутайте эти два уровня
-- команда начнёт сравнивать несопоставимые величины и~искать
потери там, где есть лишь разница между моделями учёта.

\section{Три разрыва}\label{sec:recon-gaps}

Карточная операция образует цепочку событий, и~каждая стадия
оставляет собственный след. Авторизация фиксирует право провести операцию,
клиринг закрепляет её окончательный денежный облик, а~расчёт показывает
уже фактическое движение средств с~комиссиями и взаимными вычетами.
Сверка начинается с~признания того, что эти следы по определению не обязаны
быть идентичными.

Первый разрыв возникает между авторизацией и клирингом. Ресторан может
авторизовать счёт на одну сумму, а~после чаевых отправить в~клиринг
другую; отель сначала блокирует предварительный депозит, а~потом
предъявляет итоговую стоимость проживания; авиакомпания доначисляет
часть сборов уже после исходной покупки. Поэтому сверка на этом уровне
начинается с~поиска правильной
пары: для каждой клиринговой записи надо найти соответствующую
авторизацию и проверить, укладывается ли расхождение в~допустимые
пределы~\cite{visa-tadg,mc-tpr}.

Второй разрыв возникает между клирингом и расчётной позицией. За день
может пройти сто транзакций, но в~расчётном отчёте эквайер увидит
свёрнутую итоговую сумму. Здесь многие команды впервые чувствуют, что
\enquote{деньги не сходятся}, хотя на самом деле сравнивают поток
операций с~результатом его нетто-свёртки. В~публичных правилах Mastercard нетто-расчёт описан как
базовый режим, а~прямой двусторонний расчёт между участниками вынесен
в~отдельное
исключение; Visa, со своей стороны, отдельно описывает VisaNet
Settlement Service как сервис, который формирует нетто-позицию расчёта
для участников, работающих через BASE II~\cite{mc-rules,visa-core-rules}.

Третий разрыв возникает между расчётной позицией и банковской выпиской.
Расчётная позиция в~системе схемы может уже считаться закрытой, а~движение по
корреспондентскому счёту пройти позже, в~другом окне дня или даже в
иной валюте после пересчёта. Здесь сверка работает уже не с
транзакциями как таковыми, а~с~сопоставлением расчётного файла и
банковской выписки по счёту урегулирования. Прежде чем искать
потери, точно назовите слой, на котором их ищете. Иначе спор
превратится в~обсуждение разных способов описать один и~тот же поток.
Mastercard Rules прямо требуют ежедневно сверять внутренние записи с
расчётными отчётами схемы и балансировать расхождения на этом уровне~\cite{mc-rules}.

Три разрыва за неделю.
День~1: ресторан авторизует 5\,000~руб., hold
создан. День~2: клиринговый файл приходит с~суммой
5\,200~руб. (чаевые 200~руб.); сверка сопоставляет запись
по ARN и~коду авторизации, но фиксирует расхождение
в 200~руб. -- первый разрыв. Тот же день: авторизация на
800~руб. в~другом магазине. День~5: клиринг на 800~руб.
так и~не пришёл; срок представления ещё не истёк, запись
висит в~очереди ожидания. День~8: срок представления для MCC
магазина (7~дней) истекает; запись перемещается в~очередь
исключений как orphan authorization -- hold снимается
автоматически. День~3: в~расчётном файле схемы нетто-позиция
за день составляет 4\,800~руб. (5\,200~руб. клиринга минус
400~руб. встречных обязательств) -- второй разрыв. День~4:
банковская выписка показывает зачисление 4\,795~руб. --
разница в 5~руб. объясняется комиссией за расчётное
обслуживание, которая не видна в~схемном файле -- третий
разрыв.

\section{BASE II и IPM}\label{sec:recon-base2-ipm}

Клиринговые файлы схем -- рабочие документы клирингового процесса,
а~не бухгалтерский отчёт; читайте их как машинный след
завершённых транзакций. У Visa этот слой в~публичных правилах
описан через BASE II, у Mastercard через IPM~\cite{visa-core-rules,mc-pfmi-disclosure}.

Visa определяет BASE II как файловый сервис VisaNet, который собирает,
клирит, рассчитывает и доставляет файлы финансовой и нефинансовой
активности между Visa и участниками. В~приложении к~публичным правилам
перечислены семейства форматов BASE II: TC 01--48
и TC 50--92 (с~подзаписями TCR0--TCR7 в~составе каждой записи), а~отдельно опубликованы материалы по VisaNet Settlement Service,
который выдаёт нетто-позицию расчёта и расчётные отчёты для участников
BASE II~\cite{visa-core-rules}. Жаргон вторичен; первично другое:
нужен нормализованный поток клиринга, который можно связать
с~авторизацией и последующим расчётом.

Mastercard, в~отличие от Visa, публично раскрывает именно термин
Integrated Product Messages (IPM) Clearing Formats как часть своих
стандартов, а Transaction Processing Rules регулярно отсылают к~полям
сообщения First Presentment/1240, описываемым в~руководстве
IPM Clearing Formats~\cite{mc-pfmi-disclosure,mc-tpr}. Для сверки это
даёт два следствия. Клиринг живёт не в~том же сообщении, что авторизация.
И~поля и~коды Visa BASE II не совпадают с~полями и~кодами Mastercard;
делать вид, что совпадают, нельзя.

Visa BASE II и Mastercard IPM сообщают по сути одно и~то же -- факт
завершённой операции с~её денежным обликом, -- но структура полей
различается. В BASE II идентификаторы сопоставления живут
в Transaction Code (TC) и~связке ARN (Acquirer Reference Number)
с RRN; в IPM/1240 аналогичную роль играют DE~31 (Acquirer Reference
Data) и DE~48 (Additional Data)~\cite{visa-core-rules,mc-tpr}.
Валюта транзакции и валюта расчёта кодируются в~разных полях, а
формат представления суммы (с~разделителем или без, с~фиксированной
точностью или плавающей) тоже отличается между схемами.

Прежде чем механизм сверки начнёт сопоставлять записи, приведите каждую
клиринговую запись к~единой внутренней модели: общий идентификатор операции,
сумма в~минорных единицах валюты, дата обработки в~одном часовом
поясе, код валюты по ISO~4217. Без этого шага даже идеальный алгоритм
матчинга будет регулярно спотыкаться на разнице форматов, а~не на
реальных расхождениях.

\begin{figure}[tbp]
  \centering
  \includefiguresvg[width=.8\linewidth]{assets/figures/ch23-reconciliation/base2-vs-ipm}
  \caption{Клиринговые форматы: Visa BASE II vs.\ Mastercard IPM.}\label{fig:base2-vs-ipm}
\end{figure}

Парсер должен собрать клиринговую запись, извлечь поля сопоставления
и передать их в~механизм сверки, не делая вид, будто Visa BASE II
и Mastercard IPM -- один и~тот же формат. Для каждой схемы нужен свой адаптер или
хотя бы аккуратный промежуточный слой. Иначе ошибка возникает ещё до
бухгалтерии, на этапе разбора файла.

\section{Нетто- и~валовой расчёт}\label{sec:recon-net-vs-gross}

Сумма транзакций почти никогда не совпадает с~тем, что пришло на
расчётный счёт: карточные системы показывают результат
нетто-свёртки, а~не вал операций. По умолчанию работает
нетто-расчёт (net settlement); валовой расчёт (gross settlement)
остаётся исключением.

В~модели нетто-расчёта схема сводит все позиции дня в~одну итоговую
сумму. За множеством покупок, возвратов и комиссий остаётся одна
очищенная позиция, которую и~видит эквайер. Подход уменьшает
количество отдельных переводов и делает расчётную инфраструктуру дешевле и спокойнее.

Валовой расчёт устроен иначе: каждая операция или небольшая группа
операций идёт отдельно. Режим прозрачнее, потому что деньги не
растворяются в~общей дневной сумме, но платит он другую цену:
учёт становится более детальным, а дисциплина сверки -- более строгой.

Выбор режима расчёта не принадлежит ТСП. Схема определяет
базовый режим для участника, а~эквайер транслирует его в~условия
договора. Mastercard описывает нетто-расчёт как стандартный; прямой
двусторонний расчёт возможен, но требует отдельного соглашения между
участниками~\cite{mc-rules}. Visa формирует нетто-позицию через
VisaNet Settlement Service~\cite{visa-core-rules}. Для PSP или
маркетплейса это означает: расчётный режим наследуется от
эквайера, и~менять его в~одностороннем порядке нельзя.

Внутри нетто-свёртки скрыты несколько компонентов: покупки, возвраты,
interchange fee, assessments и штрафные списания. Чтобы сверка
работала, каждый компонент должен быть выделен отдельно: итоговая
сумма по выписке = покупки минус возвраты минус interchange минус
assessments минус прочие удержания. Не раскрыт хотя бы один компонент --
дельта между ожидаемой и фактической суммой становится
неопознанным расхождением, которое команда расследует вручную каждый
расчётный день.

\begin{figure}[tbp]
  \centering
  \includefiguresvg[width=.7\linewidth]{assets/figures/ch23-reconciliation/net-vs-gross}
  \caption{Нетто- и~валовой расчёт: 100 транзакций в~одну позицию.}\label{fig:net-vs-gross}
\end{figure}

Если схема работает в~нетто-режиме,
искать в~банковской выписке сумму всех покупок за день бессмысленно.
Сначала сверяют клиринговый поток с~расчётной позицией, и~только потом
расчётную позицию с~банковской выпиской.

\practicenote{%
  Не сравнивайте авторизационный журнал сразу с~банковской выпиской:
  авторизация ещё не стала клирингом, а~выписка уже отражает расчёт.
  Правильная цепочка сверки: авторизация против клиринга, клиринг
  против расчётной позиции, расчётная позиция против выписки.

}
\section{Исключения}\label{sec:recon-exceptions}

Транзакции, которые не сходятся автоматически, прятать в~общей
массе нельзя. Зрелая сверка начинается с~явной классификации: каждое несоответствие
получает понятный тип и~свой маршрут обработки -- автоматический
или ручной.

Чаще всего первым всплывает расхождение суммы. В~ресторане
авторизовали одну сумму, а~в~клиринге пришла другая; в~гостинице
предварительная блокировка не совпала с~окончательным чеком. Такая
запись не требует паники, но требует проверки допустимого отклонения.
Порог превышен -- запись штатной уже не считается~\cite{visa-core-rules}.

Другая группа исключений: клиринг без авторизации. Так бывает при
принудительном представлении, техническом сбое или в~редких сценариях,
когда авторизация прошла вне основного процесса записи. Главное здесь --
объяснение происхождения: перед нами допустимый сценарий
или признак риска.

Обратная проблема: авторизация без клиринга. Операция была разрешена,
но в~клиринг так и~не ушла. За этим стоит ошибка маршрутизации,
сбой у~продавца или истечение окна представления. Не закроете запись
вовремя -- держатель карты увидит только авторизационную
блокировку, а~продавец так и~не получит выручку.

Ещё один сценарий: дубликаты в~клиринговом файле. Если
эквайер отправляет одну и~ту же запись дважды (ошибка
повторной передачи, сбой батч-процессора), а~сеть не
отфильтровала дубль, эмитент получает два идентичных
требования. Обнаруживает проблему обычно эмитент при сверке:
два клиринга к~одной авторизации. Стандартный путь
исправления -- возвратный клиринг (reversal presentment),
который эквайер отправляет через сеть~\cite{visa-core-rules}.
Для команды сверки дубликат опаснее расхождения в~сумме:
его легко пропустить, потому что каждая запись в~отдельности
выглядит корректной.

Наконец, есть позднее представление (Late Presentment): транзакция
дошла до схемы позже применимого срока представления. У
Mastercard это создаёт риск возвратного платежа из-за истечения периода защиты
от спора, а~у Visa до 12 апреля 2024 года было закреплено отдельное
основание спора Condition 12.1 (Late Presentment); с 13 апреля 2024
Visa объединила его с Condition 11.3 \enquote{No Authorization /
Late Presentment}~\cite{mc-tpr,visa-core-rules}. Для команды сверки это
значит: одной бухгалтерской классификации недостаточно. Нужно
ещё понять, нарушено ли окно представления и не открыло ли это спор по
правилам схемы.

\section{DCC и мультивалютный расчёт}\label{sec:recon-dcc}

DCC меняет сумму на экране покупателя и -- для учёта -- саму точку
фиксации курса. В~сценарии
динамической конвертации валюты (Dynamic Currency Conversion, DCC)
клиент выбирает валюту ещё на кассе, а~правила требуют заранее
раскрыть курс, возможную наценку или комиссию и~сам факт
опциональности DCC~\cite{visa-core-rules,mc-tpr}.

Если держатель карты
выбрал DCC, выбранная валюта становится валютой транзакции, а~эквайер
передаёт исходную сумму и~валюту до конверсии отдельными полями в~сообщении.
Если DCC не выбрана, операция должна идти в~локальной валюте точки
продаж или банкомата~\cite{visa-core-rules,mc-tpr}.

Mastercard отдельно указывает,
что корпорация конвертирует операции Interchange System в~применимую
валюту расчёта, а VisaNet Settlement Service формирует нетто-позицию
для участников после схемной обработки~\cite{mc-rules,visa-core-rules}.
Часть расхождений рождается именно из сочетания валюты транзакции,
валюты расчёта сети, даты обработки и DCC-наценки.
Сверка должна уметь отделять курсовое отклонение от
неверной проводки.

Отдельная проблема -- момент конвертации. Mastercard конвертирует
операцию в~валюту расчёта по курсу на дату обработки (processing
date), а~не на дату покупки~\cite{mc-rules}. Между покупкой и
обработкой может пройти несколько дней, и~курс успевает сдвинуться.
Даже при идентичных суммах в~валюте транзакции итоговый расчёт
в~валюте эквайера будет отличаться от ожидаемого. Visa использует
аналогичный механизм через свой конвертационный сервис~\cite{visa-core-rules}.

Практический чеклист мультивалютной сверки: (1)~зафиксировать
валюту транзакции, валюту расчёта и, если применимо, валюту
биллинга; (2)~хранить дату и~курс конвертации рядом с~записью;
(3)~отделять DCC-наценку от схемной конвертации; (4)~сверять
расчётную позицию в~валюте расчёта, а~не в~валюте транзакции.
Без любого из четырёх пунктов система будет регулярно показывать
расхождения, которые на самом деле объясняются курсовой механикой.

\begin{figure}[tbp]
  \centering
  \includefiguresvg[width=.8\linewidth]{assets/figures/ch23-reconciliation/dcc-currency-flow}
  \caption{Путь валюты в DCC-транзакции.}\label{fig:dcc-currency-flow}
\end{figure}

\section{Автоматизация сверки}\label{sec:recon-automation}

Механизм сопоставления находит одну и~ту же
операцию в~разных представлениях. Он работает в~несколько проходов:
сначала строгий матч по идентификаторам и сумме, затем более мягкая
проверка по дате и~другим полям, затем очередь исключений. Ограничение
здесь не техническое, а организационное: чем слабее критерии
сопоставления, тем аккуратнее должна быть ручная верификация.

Типичная архитектура использует три прохода. Первый -- строгий:
совпадение по ARN или схемному идентификатору, точная сумма, дата
в~пределах одного дня. Второй -- мягкий: совпадение по сумме и
карте с~допуском по дате в $\pm$2 дня, учёт частичных возвратов.
Третий -- ручная очередь: всё, что не сопоставилось автоматически.
Воронка определяет здоровье процесса: если на третий проход
регулярно приходится больше 2--3\% записей, проблема не в
алгоритме, а~в~качестве входных данных или в~пропущенном сценарии.

Позднее представление (late presentment) -- частая причина ложных
расхождений. Транзакция авторизована вчера, но клиринговая запись
появляется через два-три дня. Если сверка закрывает день жёстко
по календарной дате, такая запись останется сиротой до прихода
клиринга, а~потом потребует ручного сопоставления. Более устойчивый
подход -- оставлять окно ожидания: несопоставленные авторизации
ждут клиринга в~течение допустимого срока представления схемы
(после реформы Visa Authorization Framework 13~апреля 2024~года --
5~дней для розничного retail (CP) и~MIT, 10~дней для CNP, до~30~дней
для T\&E и EAS-сценариев с~estimated authorization
indicator~\cite{visa-core-rules,mc-tpr}) и~только после истечения
окна попадают в~очередь исключений.

Три метрики показывают состояние сверки: доля автоматически
сопоставленных записей (match rate), средний возраст
несопоставленной записи и~размер очереди исключений на конец
дня. Конкретные пороги \enquote{здоровых} значений у~каждого
процессора свои и~зависят от стабильности форматов клиринговых
файлов, доли возвратов и качества клиентских реквизитов;
универсальных публичных бенчмарков по этим метрикам в~открытых
источниках нет. Сильный сигнал даёт собственный исторический
ряд: если match rate падает, причина обычно в~изменении формата
файла или в~новом сценарии, которого парсер не знает. Если
возраст несопоставленных записей растёт -- очередь перегружена
и команда не справляется с~ручным разбором.

\begin{figure}[tbp]
  \centering
  \includefiguresvg[width=.75\linewidth]{assets/figures/ch23-reconciliation/matching-passes}
  \caption{Многопроходная сверка: строгий, мягкий, ручная очередь.}\label{fig:matching-passes}
\end{figure}

Готовый сервис сверки имеет смысл там, где нет прямого доступа к
схемам или объём транзакций не оправдывает собственной разработки.
Свой механизм строят там, где есть сложный мультивалютный
процесс, собственная ERP и специфические правила нетто-расчёта.

Команда, которая не смотрит на долю сопоставленных записей, время
закрытия дня и~размер очереди исключений, не узнает, работает ли
система как надо или молча накапливает проблемы.

\practicenote{%
  Самые упрямые ложные расхождения в~сверке связаны с~датой отсечки.
  Схемный расчёт живёт по операционному времени закрытия дня, а~не по
  календарной дате. Поэтому транзакция, прошедшая поздно вечером,
  может попасть в~следующий расчётный день. Если система жёстко
  привязана к~календарной дате, фиктивные расхождения будут появляться
  регулярно, особенно в~конце недели и~перед праздниками.

}

\section{Клиринг и~сверка ПС \enquote{Мир} через НСПК}\label{sec:recon-mir}

Архитектура ОПКЦ НСПК и~жизненный цикл внутрироссийской карточной
операции -- в~главе~\ref{ch:mir}. Здесь -- следствия для сверки.

Сверка по картам \enquote{Мир} опирается не на BASE~II или IPM,
а~на собственные клиринговые файлы НСПК~\cite{nspk-mir-rules}: формат
наследует ISO~8583, но имеет свои расширения, коды причин и
идентификаторы.\footnote{Имена файлов, периодичность формирования
и~состав полей -- закрытая часть документации НСПК; перед публикацией
сверяется с~действующей редакцией интерфейсной спецификации.} Эквайеру
и~эмитенту приходится сверяться сразу на трёх границах: между
собственным авторизационным журналом и клиринговой записью НСПК
(типичный случай -- пропущенный или поздно поступивший reversal);
между клиринговой позицией ОПКЦ и~расчётным отчётом (комиссии, плата
за услуги ОПКЦ); между расчётным отчётом и~движением по корсчёту
в~Банке России.

Кобейджинговые карты \enquote{Мир--UnionPay} и \enquote{Мир--JCB}
добавляют параллельный процесс: транзакции, прошедшие за рубежом
через сеть партнёра, возвращаются в~российский контур уже как
зарубежная операция эмитента -- со своим клиринговым форматом
(UnionPay, JCB). Эмитенту кобейджинговой программы приходится
держать несколько процессов сверки одновременно.

\section{Сверка платежей по СБП}\label{sec:recon-sbp}

C2B-операции по СБП (глава~\ref{ch:sbp}) проходят по принципиально
другой схеме, чем карточные. У~СБП нет двухтактной логики
\enquote{авторизация и потом клиринг}: исполненный перевод
сразу окончателен и сразу учтён на корсчёте получателя.
Сверка по СБП поэтому меняет акцент: она проверяет связку
между платёжной сессией продавца и фактическим зачислением
на расчётный счёт, а~не сводит авторизационный поток
с~клиринговым.

Источников данных у~продавца три: собственный журнал платёжных
сессий и~заказов; реестр операций от банка-эквайера или PSP,
полученный по протоколу ОПКЦ СБП (раздел~\ref{sec:sbp-protocol});
банковская выписка по расчётному счёту продавца. Здоровая
сверка СБП сопоставляет три потока по идентификатору
\texttt{operationId} (или его аналогу в API банка-партнёра),
проверяет совпадение сумм и~по необходимости ловит частичные
возвраты как отдельную цепочку операций.

Особенности сверки СБП относительно карт. Финальность перевода
означает, что между \enquote{ожидает подтверждения} и \enquote{исполнен}
не возникает отдельного клирингового шага: расхождения, которые
в~карточном мире сглаживаются клиринговым окном, в~СБП проявляются
немедленно и~требуют разбора до закрытия дня. Возврат сопоставляется
с~исходным платежом как отдельная операция, а~не как сторно. Сетевого
возвратного платежа в~СБП нет (см.~\ref{sec:disputes-sbp}), поэтому
позиции \enquote{спорная операция} в~реестре сверки не возникает.

B2C-выплаты по реестру требуют отдельного процесса:
сверка идёт по реестру в~целом и~по каждой строке отдельно
(идемпотентность строк, повторная отправка только ошибочных
позиций, отдельный контроль статуса каждого получателя).

\practicenote{%
  Если ваш продукт принимает оплаты и~через карты, и~через
  СБП, держите две независимые системы сверки и~не пытайтесь
  свести их в~одну универсальную модель: общая модель
  обычно теряет важные различия (например, отсутствие
  возвратного платежа в~СБП или отсутствие отложенного клиринга
  в~Мир относительно карт международной эмиссии до 2022~года).
  Общим остаётся только формальный объект -- продажа,
  с~которой связан один или несколько расчётных потоков.

}
